\documentclass[11pt,a4paper]{article}

\usepackage[margin=2.5cm]{geometry}
\usepackage{amsmath,amssymb,amsthm}
\usepackage{bm}
\usepackage{booktabs}
\usepackage{xcolor}
\usepackage{hyperref}
\usepackage{enumitem}
\usepackage{mathtools}

\newtheorem{remark}{Remark}
\newtheorem{assumption}{Assumption}

\newcommand{\R}{\mathbb{R}}
\newcommand{\dom}{\Omega}
\newcommand{\pore}[1]{\Omega_{#1}}
\newcommand{\iface}{\gamma}
\newcommand{\vel}{\bm{u}}
\newcommand{\pres}{p}
\newcommand{\stress}{\bm{\sigma}}
\newcommand{\normal}{\bm{n}}
\newcommand{\trac}{\bm{\lambda}}
\newcommand{\DtN}{\mathbf{G}}
\newcommand{\Schur}{\mathbf{S}}
\newcommand{\modeset}{\mathcal{M}}
\newcommand{\loadvec}{\bm{f}}
\newcommand{\zero}{\bm{0}}

\hypersetup{colorlinks=true,linkcolor=blue,urlcolor=blue,citecolor=blue}
\setlength{\emergencystretch}{2em}

\begin{document}

\title{\textbf{Affine-DDPNM: Revised Methodology}}
\author{}
\date{August 6, 2026}
\maketitle

\begin{center}
\fbox{\begin{minipage}{0.92\textwidth}
This revision fixes the traction/load sign convention, the orientation of
all tangential modes, the assumptions needed for local invertibility and
global positive definiteness, and the interpretation of the Galerkin
enrichment statement.  The original files are not modified.
\end{minipage}}
\end{center}

\section{Model problem and weak-form convention}

We consider steady incompressible Stokes flow in a pore domain
$\dom\subset\R^3$:
\begin{align}
 -\nu\Delta\vel+\nabla\pres&=\zero &&\text{in }\dom,\\
 \nabla\cdot\vel&=0 &&\text{in }\dom,\\
 \vel&=\zero &&\text{on }\Gamma_{\rm dir},\\
 \stress(\vel,\pres)\normal&=-p_b\normal
 &&\text{on }\Gamma_{\rm in}\cup\Gamma_{\rm out},
\end{align}
where
\[
  \stress(\vel,\pres)=-\pres I+\nu\nabla\vel.
\]
The last condition is a prescribed pseudotraction, not a pressure Dirichlet
condition.  In particular, $\normal\cdot\stress\normal=-p_b$ does not imply
$\pres=p_b$ pointwise because the normal viscous term is also present.

We use
\[
 b_i(\bm v,q)=-\int_{\pore{i}}q\,\nabla\cdot\bm v\,dx,
\]
so that the local discrete system has the block form
\[
 \begin{bmatrix}A_i&B_i^\top\\B_i&0\end{bmatrix}
 \begin{bmatrix}U_i\\P_i\end{bmatrix}
 =\begin{bmatrix}f_i\\0\end{bmatrix}.
\]
For a physical outward traction $\bm t=\stress\normal$, the weak load is
$\int_{\partial\pore{i}}\bm t\cdot\bm v\,dS$.

\section{Partition and oriented interface frames}

The domain is partitioned into non-overlapping Lipschitz subdomains
$\{\pore{i}\}_{i=1}^N$.  Each internal interface $\iface_k$ is shared by
exactly two pores, denoted $(i_k,j_k)$, and is assumed planar in the theory.
Choose once and for all a global orthonormal frame
\[
 \{\normal_k,\bm t_{k,1},\bm t_{k,2}\}
\]
on $\iface_k$, with $\normal_k$ directed from $i_k$ to $j_k$.  Let
$\sigma_{i,k}=+1$ for $i=i_k$ and $\sigma_{i,k}=-1$ for $i=j_k$.  The local
outward normal is
\[
 \normal_{i,k}=\sigma_{i,k}\normal_k.
\]

To match the implementation in which the assembled normal load direction is
$-\normal_{i,k}$ and the assembled tangential load directions are
$\sigma_{i,k}\bm t_{k,\nu}$, define the oriented inward-traction directions
\begin{equation}\label{eq:oriented-directions}
 \bm d_{i,k,n}=\normal_{i,k},\qquad
 \bm d_{i,k,t_\nu}=-\sigma_{i,k}\bm t_{k,\nu},\quad \nu=1,2.
\end{equation}
Then
\[
 \bm d_{j,k,\beta}=-\bm d_{i,k,\beta}
 \qquad (\beta=n,t_1,t_2).
\]
The minus sign in the tangential definition is only an orientation choice;
changing it together with the corresponding coefficient changes no traction
space.

For the planar interface centroid $\bm c_k$, define dimensionless coordinates
\begin{equation}\label{eq:coords}
 s_k(\bm x)=\frac{(\bm x-\bm c_k)\cdot\bm t_{k,1}}{a_k},\qquad
 t_k(\bm x)=\frac{(\bm x-\bm c_k)\cdot\bm t_{k,2}}{b_k},
\end{equation}
with nonzero scales $a_k,b_k$.  Let
\[
 \varphi_0=1,\qquad \varphi_s=s_k,\qquad \varphi_t=t_k,
 \qquad
 \modeset=\{0,s,t\}\times\{n,t_1,t_2\}.
\]

\section{Affine traction Ansatz and code-compatible load sign}

The unknown coefficient $\lambda_{k,\alpha,\beta}$ parameterizes the
\emph{inward} traction $-\stress\normal_{i,k}$:
\begin{equation}\label{eq:traction-ansatz}
 -\stress(\vel_i,\pres_i)\normal_{i,k}
 =\sum_{(\alpha,\beta)\in\modeset}
 \lambda_{k,\alpha,\beta}\,\varphi_\alpha\bm d_{i,k,\beta}.
\end{equation}
Consequently the physical outward traction entering the weak form is
\[
 \stress\normal_{i,k}
 =-\sum_{(\alpha,\beta)\in\modeset}
 \lambda_{k,\alpha,\beta}\,\varphi_\alpha\bm d_{i,k,\beta}.
\]
The primitive load vector is therefore
\begin{equation}\label{eq:primitive-load}
 \bigl[\loadvec_{i,k}^{(\alpha,\beta)}\bigr]_j
 =-\int_{\iface_k}\varphi_\alpha\,
 \bm d_{i,k,\beta}\cdot\bm\phi_j^{(i)}\,dS.
\end{equation}
This leading minus sign is essential.  For the normal mode it gives the
assembled direction $-\normal_{i,k}$; for the tangential modes, using
\eqref{eq:oriented-directions}, it gives
$+\sigma_{i,k}\bm t_{k,\nu}$, exactly the convention used in the supplied
basis implementation.

At an inlet or outlet port only the constant normal mode is retained.  The
known coefficient is $\lambda^k=p_b$, so that the physical outward traction
is $-p_b\normal$.

\section{Local finite-element response and compliance map}

On pore $\pore{i}$ use an inf-sup stable velocity-pressure pair
$(V_{i,h},Q_{i,h})$.  The rigorous finite-dimensional assumptions are:

\begin{assumption}\label{ass:local}
For every pore $i$:
\begin{enumerate}[nosep,leftmargin=*]
 \item $A_i$ is symmetric positive definite.  A positive-measure no-slip
 wall segment is a standard sufficient geometric condition for this
 coercivity.
 \item $B_i$ has full row rank in the pressure space actually used in the
 computation.  This is a discrete inf-sup/rank assumption and is not implied
 by the wall condition alone.
 \item The constant function $1$ belongs to $Q_{i,h}$.
\end{enumerate}
\end{assumption}

For each primitive mode solve
\begin{equation}\label{eq:unit-response}
 \begin{bmatrix}A_i&B_i^\top\\B_i&0\end{bmatrix}
 \begin{bmatrix}U_{i,k}^{(\alpha,\beta)}\\
 X_{i,k}^{(\alpha,\beta)}\end{bmatrix}
 =\begin{bmatrix}\loadvec_{i,k}^{(\alpha,\beta)}\\0\end{bmatrix}.
\end{equation}
Let $L_i$ collect all primitive load vectors, including one constant normal
column for each known boundary port.  Define
\begin{equation}\label{eq:H}
 H_i=A_i^{-1}-A_i^{-1}B_i^\top
 (B_iA_i^{-1}B_i^\top)^{-1}B_iA_i^{-1}.
\end{equation}
Then $U_i=H_iL_i\lambda_i$, and the local compliance matrix is
\begin{equation}\label{eq:G}
 \DtN_i=L_i^\top H_iL_i.
\end{equation}
It is symmetric positive semidefinite.  For any local coefficient vector,
\begin{equation}\label{eq:energy}
 \lambda_i^\top\DtN_i\lambda_i=U_i^\top A_iU_i\ge0.
\end{equation}

For a local mode define the weighted velocity moment
\begin{equation}\label{eq:moment}
 M_{i,k}^{(\alpha,\beta)}(U_i)
 =\int_{\iface_k}\varphi_\alpha\,
 U_{i,h}\cdot\bm d_{i,k,\beta}\,dS.
\end{equation}
Because of the load convention,
\begin{equation}\label{eq:G-moment}
 \bigl[\DtN_i\lambda_i\bigr]_{k,\alpha,\beta}
 =-M_{i,k}^{(\alpha,\beta)}(U_i).
\end{equation}
Thus $\DtN_i$ is the negative moment response associated with inward-traction
coefficients.  The sign has no effect on symmetry or energy positivity, but
it must be retained in the interpretation and in the full-trace restriction.

\section{Global Schur system}

Let $\lambda\in\R^{9m}$ collect the shared unknown coefficient blocks on all
internal interfaces.  Since $\bm d_{j,k,\beta}=-\bm d_{i,k,\beta}$, using the
same coefficient block on the two sides produces equal and opposite physical
tractions.  Continuity is imposed weakly through all nine moments:
\begin{equation}\label{eq:moment-continuity}
 M_{i,k}^{(\alpha,\beta)}+M_{j,k}^{(\alpha,\beta)}=0.
\end{equation}
Using \eqref{eq:G-moment}, Boolean assembly gives
\begin{equation}\label{eq:schur}
 \Schur\lambda=F,\qquad
 \Schur=\sum_{i=1}^N(R_i^u)^\top\DtN_i^{uu}R_i^u,
 \qquad
 F=-\sum_{i=1}^N(R_i^u)^\top\DtN_i^{uk}R_i^k\lambda^k.
\end{equation}

\begin{remark}[What is actually required for SPD]\label{rem:spd}
Per-face linear independence of the nine scalar/vector functions is not
sufficient for positive definiteness.  The required local algebraic condition
is
\begin{equation}\label{eq:rank-condition}
 \ker(H_iL_i)=\operatorname{span}\{\bm 1_{i,0n}\},
\end{equation}
where $\bm 1_{i,0n}$ has value one in every constant-normal port slot and
zero in all other slots.  Equivalently,
$\operatorname{rank}(H_iL_i)=n_i-1$.  This condition excludes dependencies
between different ports and dependencies created after the incompressible
projection $H_i$.  It must be proved for a specific discretization or checked
numerically; it does not follow merely from independence on each face.
\end{remark}

Under Assumption~\ref{ass:local}, condition~\eqref{eq:rank-condition}, and the
graph anchoring condition that every connected pore component contains at
least one known-traction port, the companion mathematical document proves
that $\Schur$ is SPD.  The proof uses zero extension into the known-port slots;
it does not apply the full local-kernel result directly to a shorter unknown
vector.

Floating basins and graph anchoring are distinct issues.  Merging a basin with
no wall facet is intended to restore local coercivity of $A_i$.  Ensuring that
each connected pore component touches an inlet or outlet removes the remaining
global constant-pressure gauge.  One operation does not logically imply the
other.

\section{Algorithm}

\begin{enumerate}[leftmargin=*,itemsep=2pt]
 \item For each pore, assemble $A_i$, $B_i$, and verify the intended pressure
 rank convention.
 \item For every internal port and every $(\alpha,\beta)\in\modeset$, assemble
 the signed load \eqref{eq:primitive-load} and solve
 \eqref{eq:unit-response}.  For each known boundary port solve the single
 constant-normal response.
 \item Form $\DtN_i=L_i^\top H_iL_i$.  As a diagnostic, inspect the singular
 values of $H_iL_i$ or eigenvalues of $\DtN_i$ and verify that exactly one
 local null direction is present when \eqref{eq:rank-condition} is assumed.
 \item Assemble \eqref{eq:schur}.  Check symmetry and the smallest eigenvalue
 or a Cholesky factorization before treating the system as SPD.
 \item Solve for $\lambda$ and reconstruct each local field by linear
 superposition of primitive responses.
 \item Verify local mass residuals, interface moment residuals, and the global
 boundary flux balance.
\end{enumerate}

\section{Relation to classic DD-PNM and scope of enrichment}

Let $E_{\rm P0}:\R^m\to\R^{9m}$ insert each classic coefficient into the
constant-normal slot.  Restricting both trial and test spaces gives
\[
 \Schur_{\rm P0}=E_{\rm P0}^\top\Schur E_{\rm P0},\qquad
 F_{\rm P0}=E_{\rm P0}^\top F.
\]
The full affine solution generally does not have the classic solution as its
constant-normal subvector.

The valid enrichment statement is conditional and reference-dependent.  If a
well-defined full interface-traction Schur system $\widehat S\xi=\widehat F$
is introduced with the same local pressure spaces and sign conventions, then
the affine system is its Galerkin restriction and the affine error is no larger
than the classic error in the $\widehat S$ energy \emph{seminorm}.  This does
not imply monotonic improvement of velocity $L^2$, pressure $L^2$, outlet
flux, or broken $H^1$ errors.

No theorem in the present analysis establishes equality between that
broken-pressure traction system and the monolithic continuous-pressure
Taylor--Hood system.  The supplied ``Exact FE-trace Schur'' computation is a
primal Schur reduction of the monolithic system and therefore verifies its own
algebraic elimination, not the missing equivalence with the broken-pressure
traction formulation.

\section{Computational cost}

A pore with $m_i^u$ internal ports and $m_i^k$ known boundary ports has
$n_i=9m_i^u+m_i^k$ primitive columns.  The offline response count therefore
grows by approximately a factor of nine relative to a one-mode internal
interface model.  The global unknown count is $9m$.  Sparse factorization cost
depends on graph fill and cannot in general be inferred from the unknown-count
factor alone.

\begin{thebibliography}{9}
\bibitem{sun2025ddpnm}
S.~Sun, Z.~Wang, L.~Zhang, and J.~Zhao.
\newblock A domain decomposition approach to pore-network modeling of porous
media flow.
\newblock \emph{arXiv:2510.13429v2}, 2026.
\end{thebibliography}

\end{document}
